\subsection*{Tail notation and the meaning of infinite bounds}
For a real sequence $(x_n)$ and an integer $N\ge1$, write
\[
 X_N=\{x_n:n>N\},\qquad
 u_N(x)=\inf X_N,\qquad v_N(x)=\sup X_N.
\]
The set $X_N$ is nonempty because it contains $x_{N+1}$. We use the
extended-bound conventions of Section~5: $u_N(x)=-\infty$ if the
tail is unbounded below, and $v_N(x)=+\infty$ if it is unbounded
above. For a nonempty set of real terms, a tail infimum cannot be
$+\infty$, and a tail supremum cannot be $-\infty$.
By Definition~10.6,
\[
 \liminf_{n\to\infty}x_n=\lim_{N\to\infty}u_N(x),\qquad
 \limsup_{n\to\infty}x_n=\lim_{N\to\infty}v_N(x).
\]
Removing more terms enlarges the infimum and decreases the supremum,
since $X_{N+1}\subseteq X_N$.

We will also use a finite-prefix observation. If a tail of a real
sequence has an upper bound $U$, then adjoining its finitely many
omitted terms still gives a bounded-above set: take the maximum of
$U$ and those finitely many terms. The corresponding minimum proves
the analogous statement for lower bounds. Thus a sequence that is
unbounded above, or unbounded below, has the same property on every
tail. Infinite symbols will be handled through order and these
conventions, rather than through real-number algebra.

\exercise{12.1}{Eventual comparison passes to both limiting bounds}

\goal Suppose there is an $N_0$ such that $s_n\le t_n$ for all
$n>N_0$. Prove
\[
 \liminf s_n\le\liminf t_n,
 \qquad
 \limsup s_n\le\limsup t_n.
\]
No boundedness assumption is imposed on the two sequences.

\step{Step 1: make the eventual comparison valid on the tails under discussion.}
Enlarge $N_0$, if necessary, to a positive integer. This keeps the
hypothesis true, because it discards additional initial terms.
Fix any integer $N\ge N_0$. For every $n>N$ we then have
$s_n\le t_n$.

For the upper tail bounds, if $v_N(t)=+\infty$, the extended-order
comparison $v_N(s)\le v_N(t)$ holds immediately. Otherwise $v_N(t)$
is real, and
\[
 s_n\le t_n\le v_N(t)\qquad(n>N).
\]
Thus $v_N(t)$ is an upper bound of the tail of $(s_n)$, so its least
upper bound satisfies $v_N(s)\le v_N(t)$.

For the lower tail bounds, if $u_N(s)=-\infty$, the comparison
$u_N(s)\le u_N(t)$ is automatic in the extended order.
Otherwise $u_N(s)$ is real, and
\[
 u_N(s)\le s_n\le t_n\qquad(n>N).
\]
Thus $u_N(s)$ is a lower bound of the tail of $(t_n)$, so it is at
most that tail's greatest lower bound. We have proved
\[
 u_N(s)\le u_N(t),\qquad v_N(s)\le v_N(t)
 \qquad(N\ge N_0).
\]

\step{Step 2: pass to the lower limiting bounds, checking infinite cases.}
If $(s_n)$ is unbounded below, every tail infimum $u_N(s)$ equals
$-\infty$. Therefore $\liminf s_n=-\infty\le\liminf t_n$, and
there is nothing further to prove for the lower bounds in this case.

If $(s_n)$ is bounded below, choose a real lower bound $a$.
For $n>N_0$, the hypothesis gives $t_n\ge s_n\ge a$. The
finite-prefix observation shows that $(t_n)$ is bounded below too.
Consequently both $u_N(s)$ and $u_N(t)$ are real for every $N$.
Their limits exist in the extended sense by Definition~10.6.
Apply Exercise~9.9(c), the order rule for existing limits of
real-valued sequences, to $u_N(s)\le u_N(t)$:
\[
 \lim_{N\to\infty}u_N(s)\le\lim_{N\to\infty}u_N(t).
\]
The limits may be $+\infty$, but the terms to which the exercise
was applied are real. By definition this is
$\liminf s_n\le\liminf t_n$.

\step{Step 3: pass to the upper limiting bounds, checking infinite cases.}
If $(t_n)$ is unbounded above, then every $v_N(t)=+\infty$, so
$\limsup s_n\le+\infty=\limsup t_n$.
If $(t_n)$ is bounded above by a real number $b$, then
$s_n\le t_n\le b$ for $n>N_0$. The finite-prefix observation
shows that $(s_n)$ is bounded above as well. Hence $v_N(s)$ and
$v_N(t)$ are real-valued sequences, with existing extended limits.
Applying Exercise~9.9(c) to their eventual inequality gives
\[
 \lim_{N\to\infty}v_N(s)\le\lim_{N\to\infty}v_N(t),
\]
which is $\limsup s_n\le\limsup t_n$.

\answer{Both comparisons hold:
\[
 \boxed{\liminf s_n\le\liminf t_n,
 \qquad\limsup s_n\le\limsup t_n}.
\]
The proof includes infinite limiting bounds. It does not treat
$+\infty$ or $-\infty$ as individual real terms in an application
of a real-sequence limit rule.}

\exercise{12.2}{A zero limit superior of absolute values}

\goal Prove the equivalence
\[
 \limsup\abs{s_n}=0\quad\Longleftrightarrow\quad s_n\longrightarrow0.
\]

\step{Forward direction, Step 1: derive the boundedness needed for finite tail suprema.}
Assume $\limsup\abs{s_n}=0$ and put $a_n=\abs{s_n}\ge0$.
If $(a_n)$ were unbounded above, every tail would be unbounded above
by the finite-prefix observation. Its every tail supremum would
therefore be $+\infty$, giving $\limsup a_n=+\infty$, contrary to
the assumed value $0$. Thus $(a_n)$ is bounded above.
Because all $a_n\ge0$, it is bounded below as well.

It follows that the numbers
\[
 V_N=\sup\{\abs{s_n}:n>N\}
\]
are finite and nonnegative for every positive integer $N$.
Their nonnegativity follows, for example, by comparing $V_N$ with
the member $\abs{s_{N+1}}\ge0$ of the tail.
The hypothesis and Definition~10.6 now say $V_N\to0$.

\step{Forward direction, Step 2: control every term in one sufficiently late tail.}
Let $\varepsilon>0$. From $V_N\to0$, choose a positive integer $N_*$
large enough that $\abs{V_{N_*}-0}<\varepsilon$.
Since $V_{N_*}\ge0$, this means $V_{N_*}<\varepsilon$.
For every integer $n>N_*$, the definition of tail supremum gives
\[
 \abs{s_n-0}=\abs{s_n}\le V_{N_*}<\varepsilon.
\]
This is Definition~7.1 for $s_n\to0$. Boundedness was obtained from
the hypothesis before any finite-tail estimate was used.

\step{Reverse direction, Step 1: verify convergence of the absolute values.}
Assume $s_n\to0$. Let $\varepsilon>0$. The definition gives a
threshold $N$ such that $n>N$ implies $\abs{s_n-0}<\varepsilon$.
Since $\abs{s_n}\ge0$, the error in the sequence of absolute values
is exactly
\[
 \abs{\abs{s_n}-0}=\abs{s_n}=\abs{s_n-0}<\varepsilon.
\]
Thus $\abs{s_n}\to0$ directly from the definition.

\step{Reverse direction, Step 2: identify its limit superior.}
Theorem~10.7 says that when a sequence has a limit, its limit
inferior and limit superior equal that limit. Applying it to the
convergent sequence $(\abs{s_n})$ gives $\limsup\abs{s_n}=0$.

\answer{The two implications prove
\(\boxed{\limsup\abs{s_n}=0\ \Longleftrightarrow\ s_n\to0}\).}

\exercise{12.3}{Two synchronized sequences of period four}

\goal Compute all seven requested combinations for the sequences
whose repeated blocks are $(0,1,2,1)$ and $(2,1,1,0)$.

\step{Step 1: keep the phases aligned in one residue table.}
The index $n$ has one of four residues modulo $4$. Reading the
entries at the same index gives
\begin{center}
\begin{tabular}{@{}ccccc@{}}
\toprule
Residue of $n$ & $s_n$ & $t_n$ & $s_n+t_n$ & $s_nt_n$\\
\midrule
$1$ & $0$ & $2$ & $2$ & $0$\\
$2$ & $1$ & $1$ & $2$ & $1$\\
$3$ & $2$ & $1$ & $3$ & $2$\\
$0$ & $1$ & $0$ & $1$ & $0$\\
\bottomrule
\end{tabular}
\end{center}
For instance, at indices $n=4j+3$ the entries are $s_n=2$ and
$t_n=1$, so the sum is $3$ and the product is $2$.
Every row occurs at arbitrarily large indices: by taking the integer
$j$ sufficiently large, its index exceeds any prescribed threshold.
For the residue $0$, use indices $4,8,12,\ldots$.

\step{Step 2: identify each tail's minimum and maximum.}
Every tail of $(s_n)$ contains the values $0,1,2$ and no others.
Thus its infimum is $0$ and its supremum is $2$. The same is true
for every tail of $(t_n)$. The sum has the repeated block
$(2,2,3,1)$, so every tail of the sum contains exactly the values
$1,2,3$; its infimum is $1$ and its supremum is $3$.
The product has the repeated block $(0,1,2,0)$, so its tail infima
are $0$ and its tail suprema are $2$.
Taking the limits of these constant sequences of tail bounds gives
\[
 \liminf s_n=0,\quad\limsup s_n=2,\qquad
 \liminf t_n=0,\quad\limsup t_n=2,
\]
\[
 \liminf(s_n+t_n)=1,\quad\limsup(s_n+t_n)=3,
 \qquad
 \liminf(s_nt_n)=0,\quad\limsup(s_nt_n)=2.
\]

\step{Step 3: evaluate every requested expression.}
\begin{align*}
 \text{(a)}\quad&\liminf s_n+\liminf t_n=0+0=\boxed{0},\\
 \text{(b)}\quad&\liminf(s_n+t_n)=\boxed{1},\\
 \text{(c)}\quad&\liminf s_n+\limsup t_n=0+2=\boxed{2},\\
 \text{(d)}\quad&\limsup(s_n+t_n)=\boxed{3},\\
 \text{(e)}\quad&\limsup s_n+\limsup t_n=2+2=\boxed{4},\\
 \text{(f)}\quad&\liminf(s_nt_n)=\boxed{0},\\
 \text{(g)}\quad&\limsup(s_nt_n)=\boxed{2}.
\end{align*}

\answer{The phase alignment matters. The largest value of $(s_n)$
occurs when $t_n=1$, whereas the largest value of $(t_n)$ occurs
when $s_n=0$. Adding the separate maxima therefore gives $4$, even
though the paired sum has maximum $3$.}

\exercise{12.4}{The limit superior of a sum is at most the sum of the limit superiors}

\goal For bounded real sequences $(s_n)$ and $(t_n)$, prove
\[
 \limsup(s_n+t_n)\le\limsup s_n+\limsup t_n.
\]

\step{Step 1: establish that all relevant tail bounds are finite.}
Choose real constants $A,B\ge0$ with $\abs{s_n}\le A$ and
$\abs{t_n}\le B$ for every $n$. The triangle inequality gives
$\abs{s_n+t_n}\le A+B$, so the sum sequence is bounded as well.
For every positive integer $N$, the three tails are nonempty and
bounded, and hence the numbers
\[
 S_N=\sup\{s_n:n>N\},\quad T_N=\sup\{t_n:n>N\},\quad
 W_N=\sup\{s_n+t_n:n>N\}
\]
are real. Their limits as $N\to\infty$ are also finite, since the
underlying sequences are bounded.

\step{Step 2: prove the finite-tail inequality.}
For every integer $n>N$, the upper-bound properties give
$s_n\le S_N$ and $t_n\le T_N$. Adding yields
\[
 s_n+t_n\le S_N+T_N.
\]
Thus the real number $S_N+T_N$ is an upper bound of the paired-sum
tail. Its least upper bound satisfies
\[
 W_N\le S_N+T_N.
\]
This argument does not require $S_N$ or $T_N$ to be attained. It
also does not require a single index where both separate bounds
are approached at once.

\step{Step 3: take limits of the real tail-bound sequences.}
By Definition~10.6,
\[
 S_N\to\limsup s_n,\qquad T_N\to\limsup t_n,
 \qquad W_N\to\limsup(s_n+t_n).
\]
The sum law from Section~9 gives
\[
 S_N+T_N\longrightarrow\limsup s_n+\limsup t_n.
\]
Apply Exercise~9.9(c) to the termwise comparison $W_N\le S_N+T_N$.
The resulting comparison of limits is
\[
 \limsup(s_n+t_n)\le\limsup s_n+\limsup t_n.
\]

\answer{This proves the desired inequality. The boundedness
hypotheses ensure that every sum of limiting bounds used here is a
sum of real numbers; no undefined combination of infinite symbols
has been introduced.}

\exercise{12.5}{The corresponding lower bound for the limit inferior of a sum}

\goal Use Exercises~11.8 and 12.4 to prove, for bounded real
sequences,
\[
 \liminf(s_n+t_n)\ge\liminf s_n+\liminf t_n.
\]

\step{Step 1: verify that the reflected sequences meet the hypotheses.}
The sequences $(-s_n)$ and $(-t_n)$ are bounded because
$\abs{-s_n}=\abs{s_n}$ and $\abs{-t_n}=\abs{t_n}$.
Their sum is $-s_n-t_n=-(s_n+t_n)$, which is bounded too.
Exercise~12.4 therefore applies and gives
\[
 \limsup\bigl(-(s_n+t_n)\bigr)
    \le\limsup(-s_n)+\limsup(-t_n).
\]

\step{Step 2: replace each reflected upper limiting bound.}
Exercise~11.8 states
$\liminf x_n=-\limsup(-x_n)$.
For the bounded sequences here, all these values are real. Taking
their negatives also gives
$\limsup(-x_n)=-\liminf x_n$. Substitute this identity in the
three places in Step~1:
\[
 -\liminf(s_n+t_n)
    \le-\liminf s_n-\liminf t_n
    =-\bigl(\liminf s_n+\liminf t_n\bigr).
\]

\step{Step 3: reverse the inequality by negation.}
Multiplying a real inequality by $-1$ reverses its direction.
Negating both sides of the last comparison, and taking the negative
twice in each expression, yields
\[
 \liminf(s_n+t_n)\ge\liminf s_n+\liminf t_n.
\]

\answer{The claimed lower inequality follows. The sign reversal is
the reason that reflecting the upper inequality from Exercise~12.4
produces a lower inequality here.}

\exercise{12.6}{Scaling the upper and lower limiting bounds}

\goal For bounded $(s_n)$ and $k\ge0$, prove (a)
$\limsup(ks_n)=k\limsup s_n$ and (b)
$\liminf(ks_n)=k\liminf s_n$. Determine the correct formulas in
(c) when $k<0$.

\step{Preliminary step: scaling preserves boundedness.}
Choose $A\ge0$ with $\abs{s_n}\le A$ for every $n$.
For any fixed real $k$, the identity
$\abs{ks_n}=\abs{k}\abs{s_n}\le\abs{k}A$ shows that $(ks_n)$
is bounded. Its tails are nonempty. Hence all tail bounds for
$(s_n)$ and $(ks_n)$ below are real numbers.
Write
\[
 m_N=\inf\{s_n:n>N\},\qquad M_N=\sup\{s_n:n>N\}.
\]
The limits of these two sequences are $\liminf s_n$ and
$\limsup s_n$, respectively, and are finite by boundedness.

\step{Parts (a) and (b), the case $k=0$.}
If $k=0$, every scaled term is zero. Every scaled tail is the
singleton $\{0\}$, whose infimum and supremum are zero.
Therefore both limiting bounds of $(ks_n)$ are zero. The expressions
$0\cdot\liminf s_n$ and $0\cdot\limsup s_n$ are also zero,
because their other factors are finite real numbers. Thus both
requested formulas hold in this case.

\step{Part (a), Step 1: prove the upper-tail scaling formula for $k>0$.}
Fix $k>0$ and an integer $N\ge1$. Since $s_n\le M_N$ for all
$n>N$, multiplication by $k$ gives $ks_n\le kM_N$.
Thus $kM_N$ is an upper bound of the scaled tail.

To check that it is the least upper bound, let $U$ be any real upper
bound of that tail. Then $ks_n\le U$ for every $n>N$.
Division by $k>0$ gives $s_n\le U/k$ for all these indices, so
$U/k$ is an upper bound of the original tail. The least-upper-bound
property gives $M_N\le U/k$. Multiplying by $k>0$ gives
$kM_N\le U$. We have proved both defining properties of a
supremum, so
\[
 \sup\{ks_n:n>N\}=kM_N.
\]

\step{Part (a), Step 2: take the limit of the real tail bounds.}
By the constant-multiple limit law from Section~9,
\[
 kM_N\longrightarrow k\limsup s_n.
\]
Definition~10.6 identifies the limit of the left side of the tail
identity as $\limsup(ks_n)$. Hence
\[
 \boxed{\limsup(ks_n)=k\limsup s_n}\qquad(k>0).
\]
Together with the zero case, this proves part (a) for all $k\ge0$.

\step{Part (b), Step 1: prove the lower-tail scaling formula for $k>0$.}
For every $n>N$, $m_N\le s_n$; multiplying by $k>0$ gives
$km_N\le ks_n$. Thus $km_N$ is a lower bound of the scaled tail.
If $L$ is any real lower bound of that tail, then $L\le ks_n$
for every $n>N$. Divide by $k>0$ to obtain $L/k\le s_n$.
Consequently $L/k$ is a lower bound of the original tail. Since
$m_N$ is its greatest lower bound, $L/k\le m_N$, and hence
$L\le km_N$. These two properties establish
\[
 \inf\{ks_n:n>N\}=km_N.
\]
Taking limits of this equality and using the constant-multiple law
gives
\[
 \boxed{\liminf(ks_n)=k\liminf s_n}\qquad(k>0).
\]
The zero case already proved extends it to all $k\ge0$.

\step{Part (b), Step 2: relate the formula to the reflection hint.}
For completeness, the same lower formula follows from part (a)
and Exercise~11.8. When $k>0$,
\begin{align*}
 \liminf(ks_n)
   &=-\limsup(-ks_n),\\
   &=-\limsup\bigl(k(-s_n)\bigr),\\
   &=-k\limsup(-s_n),\\
   &=k\liminf s_n.
\end{align*}
The third equality uses part (a) for the bounded sequence $(-s_n)$;
the first and last use reflection. All negatives here are negatives
of finite real numbers.

\step{Part (c), Step 1: a negative multiplier exchanges the tail bounds.}
Now let $k<0$ and fix $N$. Multiplication by $k$ reverses order.
Since $m_N\le s_n$ for every $n>N$, we have $ks_n\le km_N$;
thus $km_N$ is an upper bound of the scaled tail.
For any other upper bound $U$, the inequalities $ks_n\le U$
become $s_n\ge U/k$ after division by $k<0$.
Hence $U/k$ is a lower bound of the original tail, so
$U/k\le m_N$. Multiplying this last inequality by $k<0$ reverses
it and gives $U\ge km_N$. Therefore
\[
 \sup\{ks_n:n>N\}=km_N\qquad(k<0).
\]

Similarly, $s_n\le M_N$ gives $ks_n\ge kM_N$, so $kM_N$ is a
lower bound of the scaled tail. If $L$ is any other lower bound,
$L\le ks_n$ gives $L/k\ge s_n$ after division by $k<0$.
Thus $L/k$ is an upper bound of the original tail and
$M_N\le L/k$. Multiplying by $k<0$ gives $kM_N\ge L$.
It follows that
\[
 \inf\{ks_n:n>N\}=kM_N\qquad(k<0).
\]

\step{Part (c), Step 2: take finite limits.}
The constant-multiple law, now with a fixed negative constant, still
applies to these real convergent tail-bound sequences. Taking limits
of the two identities gives
\[
 \boxed{\limsup(ks_n)=k\liminf s_n,
 \qquad\liminf(ks_n)=k\limsup s_n}\qquad(k<0).
\]

\answer{A positive multiplier preserves the upper and lower limiting
bounds separately. A zero multiplier makes both zero. A negative
multiplier exchanges their roles, as the two formulas in part (c)
state. All statements were proved without assuming that any tail
supremum or infimum is attained.}

\exercise{12.7}{A positive multiplier preserves an infinite limit superior}

\goal Suppose $\limsup s_n=+\infty$ and $k>0$. Prove
$\limsup(ks_n)=+\infty$ without assuming that $(s_n)$ is bounded.

\step{Step 1: prove that every tail of the original sequence is unbounded above.}
Suppose some tail were bounded above by a real $U$. Then $s_n\le U$
for every index in that tail, so the eventual comparison result of
Exercise~12.1, applied to $(s_n)$ and the constant sequence $U$,
would give
\[
 \limsup s_n\le\limsup U=U.
\]
This contradicts $\limsup s_n=+\infty$. Thus no such tail bound
exists: for each integer $N$ and each real $R$, some index $n>N$
satisfies $s_n>R$.

\step{Step 2: prove that every scaled tail is unbounded above.}
Fix any positive integer $N$ and any proposed real upper bound $M$
for the scaled tail. Since $k>0$, the quotient $M/k$ is real.
Step~1, with $R=M/k$, gives an index $n>N$ such that $s_n>M/k$.
Multiplying by the positive number $k$ preserves strict inequality:
\[
 ks_n>M.
\]
Thus $M$ is not an upper bound. Because $M$ was arbitrary, the
scaled tail is unbounded above. Because $N$ was arbitrary, every
scaled tail has this property.

\step{Step 3: apply the definition of the infinite tail bound.}
For every $N$, the extended-bound convention gives
\[
 \sup\{ks_n:n>N\}=+\infty.
\]
Its limiting value is therefore $+\infty$, which is exactly
$\limsup(ks_n)$ by Definition~10.6.

\answer{We obtain \(\boxed{\limsup(ks_n)=+\infty}\).
The proof divided and multiplied only real numbers, such as $M$ and
$M/k$, and used unboundedness to interpret the infinite result.}

\exercise{12.8}{A product bound for bounded nonnegative sequences}

\goal Let $(s_n)$ and $(t_n)$ be bounded sequences with $s_n\ge0$
and $t_n\ge0$ for every $n$. Prove
\[
 \limsup(s_nt_n)\le(\limsup s_n)(\limsup t_n).
\]

\step{Step 1: give finite bounds for the three sequences.}
Choose $A,B\ge0$ such that $0\le s_n\le A$ and $0\le t_n\le B$
for every $n$. Because $t_n\ge0$, multiplication of $s_n\le A$
by $t_n$ gives $s_nt_n\le At_n$. Because $A\ge0$, multiplication
of $t_n\le B$ by $A$ gives $At_n\le AB$. Thus
\[
 0\le s_nt_n\le At_n\le AB.
\]
The product sequence is bounded and nonnegative too.
For every $N$, all three tails are nonempty and bounded, and their
suprema
\[
 S_N=\sup\{s_n:n>N\},\quad
 T_N=\sup\{t_n:n>N\},\quad
 P_N=\sup\{s_nt_n:n>N\}
\]
are finite and nonnegative. For example,
$0\le s_{N+1}\le S_N$, which proves $S_N\ge0$; the other two
nonnegativity statements follow in the same way from their tails.

\step{Step 2: prove the paired-product tail inequality.}
For each $n>N$, the tail upper bounds give $s_n\le S_N$ and
$t_n\le T_N$. Multiply the first inequality by $t_n\ge0$:
\[
 s_nt_n\le S_Nt_n.
\]
Multiply the second by $S_N\ge0$:
\[
 S_Nt_n\le S_NT_N.
\]
The sign checks are necessary for both comparisons. Combining them
shows $s_nt_n\le S_NT_N$ for every $n>N$.
Hence $S_NT_N$ is an upper bound of the product tail. Its least
upper bound satisfies
\[
 P_N\le S_NT_N.
\]
Neither separate tail supremum needs to be attained, and no common
index attaining both was selected.

\step{Step 3: compare the finite limits of the tail-bound sequences.}
Definition~10.6 gives
\[
 S_N\to\limsup s_n,\qquad T_N\to\limsup t_n,
 \qquad P_N\to\limsup(s_nt_n).
\]
All limits here are finite because the original sequences and the
product sequence are bounded. The product limit law from Section~9
therefore yields
\[
 S_NT_N\longrightarrow(\limsup s_n)(\limsup t_n).
\]
Applying Exercise~9.9(c) to $P_N\le S_NT_N$ gives the requested
comparison of limiting values.

\answer{We have proved
\[
 \boxed{\limsup(s_nt_n)\le(\limsup s_n)(\limsup t_n)}.
\]
The proof includes the case where either upper limiting bound is
zero, because no division by one of those bounds was used.}
